% GENERATED FILE. Edit canonical lessons, then run npm run generate:books.
% canonical-source-hash: fnv1a64:9cd6621eefb58bd6
% canonical-lessons: ZH-C273-chengshi, ZH-C273-haoke, ZH-C273-yongxin, ZH-C273-cuxin, ZH-C273-xixin

\chapter{Honest and Careful}
\label{ch:zh-honest-and-careful}
\hlchaptermodality{273}

\begin{chapteropening}
\textbf{By the end of this chapter:} \emph{I can say that someone is honest, hospitable, careless, or careful and attentive.}

Complete the last lesson of chapter 273: I can say that someone is honest, hospitable, careless, or careful and attentive.
\end{chapteropening}

\section[chéngshí]{chéngshí — honest}

\label{lesson:ZH-C273-chengshi}



\begin{quote}
Before the new one: say the Mandarin for an uncle, then the Mandarin for an auntie.
\end{quote}

\subsection*{You'll want to know: chéngshí}
\emph{chéngshí} — \textquotedblleft{}honest\textquotedblright{}.

\begin{quote}
\emph{tā shì yí ge chéngshí de rén} — He's an honest person.
\end{quote}

Its characters come later; for now it is a word for the ear.

\begin{cousinweb}
\emph{chéngshí} is built from \emph{chéng} and \emph{shí}, both inherited from Middle Chinese, \emph{dzyeng} and \emph{zyit} (Guangyun, in Baxter's transcription), which glosses the first 'trustworthy' and the second 'full' and 'sincere'.
\end{cousinweb}

\begin{grammarlens}[title={What you've built}]
One more word for describing things.
\end{grammarlens}

\subsection*{Your turn}
\begin{itemize}
\raggedright
  \item \emph{Say it:} \emph{chéngshí}
  \item \emph{Say it:} \emph{chéngshí}, once more
  \item \emph{Say it:} \emph{tā shì yí ge chéngshí de rén} — He's an honest person.
  \item \emph{Recall:} say the Mandarin for an uncle, then the Mandarin for an auntie, then say \emph{chéngshí} again
\end{itemize}

\begin{tcolorbox}[breakable,colback=teal!4,colframe=teal!35!black,title={Before you move on}]
What does chéngshí mean? (\textquotedblleft{}Honest\textquotedblright{}.) Say it once more.
\end{tcolorbox}
\section[\texorpdfstring{\zh{好客} (hàokè)}{好客 (hàokè)}]{\zh{好客} (hàokè) — hospitable}

\label{lesson:ZH-C273-haoke}



\begin{quote}
Before the new one: say the Mandarin for an auntie, then the Mandarin for honest.
\end{quote}

\subsection*{You'll want to know: \zh{好客}}
\textbf{\zh{好客}} — \emph{hàokè} — \textquotedblleft{}hospitable\textquotedblright{}.

\begin{quote}
\emph{tāmen jiā fēicháng hàokè} — Their family is very hospitable.
\end{quote}

\begin{cousinweb}
\textbf{\zh{好客}} (\emph{hàokè}) is built from \zh{好}, here said with a falling tone, 'to be fond of', and \zh{客} 'a guest', both inherited from Middle Chinese, \emph{xawH} and \emph{khaek} (Guangyun, in Baxter's transcription), which glosses the second 'a guest'.
\end{cousinweb}

\begin{grammarlens}[title={What you've built}]
One more word for describing things.
\end{grammarlens}

\subsection*{Your turn}
\begin{itemize}
\raggedright
  \item \emph{Say it:} \emph{hàokè}
  \item \emph{Say it:} \emph{hàokè}, once more
  \item \emph{Say it:} \emph{tāmen jiā fēicháng hàokè} — Their family is very hospitable.
  \item \emph{Recall:} say the Mandarin for an auntie, then the Mandarin for honest, then say \emph{hàokè} again
\end{itemize}

\begin{tcolorbox}[breakable,colback=teal!4,colframe=teal!35!black,title={Before you move on}]
What does \zh{好客} mean? (\textquotedblleft{}Hospitable\textquotedblright{}.) Say it once more.
\end{tcolorbox}
\section[\texorpdfstring{\zh{用心} (yòngxīn)}{用心 (yòngxīn)}]{\zh{用心} (yòngxīn) — attentive; putting your heart into it}

\label{lesson:ZH-C273-yongxin}



\begin{quote}
Before the new one: say the Mandarin for honest, then the Mandarin for hospitable.
\end{quote}

\subsection*{You'll want to know: \zh{用心}}
\textbf{\zh{用心}} — \emph{yòngxīn} — \textquotedblleft{}attentive; putting your heart into it\textquotedblright{}.

\begin{quote}
\textbf{\zh{我老公工作很用心}} \emph{wǒ lǎogōng gōngzuò hěn yòngxīn} — My husband puts his heart into his work.
\end{quote}

\begin{cousinweb}
\textbf{\zh{用心}} (\emph{yòngxīn}) is built from \zh{用} 'to use' and \zh{心} 'the heart', both inherited from Middle Chinese, \emph{yowngH} and \emph{sim} (Guangyun, in Baxter's transcription).
\end{cousinweb}

\begin{grammarlens}[title={What you've built}]
One more word for describing things.
\end{grammarlens}

\subsection*{Your turn}
\begin{itemize}
\raggedright
  \item \emph{Say it:} \emph{yòngxīn}
  \item \emph{Say it:} \emph{yòngxīn}, once more
  \item \emph{Say it:} \emph{wǒ lǎogōng gōngzuò hěn yòngxīn} — My husband puts his heart into his work.
  \item \emph{Recall:} say the Mandarin for honest, then the Mandarin for hospitable, then say \emph{yòngxīn} again
\end{itemize}

\begin{tcolorbox}[breakable,colback=teal!4,colframe=teal!35!black,title={Before you move on}]
What does \zh{用心} mean? (\textquotedblleft{}Attentive; putting your heart into it\textquotedblright{}.) Say it once more.
\end{tcolorbox}
\section[cūxīn]{cūxīn — careless}

\label{lesson:ZH-C273-cuxin}



\begin{quote}
Before the new one: say the Mandarin for hospitable, then the Mandarin for putting your heart into it.
\end{quote}

\subsection*{You'll want to know: cūxīn}
\emph{cūxīn} — \textquotedblleft{}careless\textquotedblright{}.

\begin{quote}
\emph{nǐ tài cūxīn le} — You're so careless.
\end{quote}

Its characters come later; for now it is a word for the ear.

\begin{cousinweb}
\emph{cūxīn} is built from \emph{cū} 'coarse' and \zh{心} 'the heart', both inherited from Middle Chinese, \emph{dzuX} and \emph{sim} (Guangyun, in Baxter's transcription), which glosses the first 'rough'. Middle Chinese \emph{dzuX} would regularly give zù; the initial and the level tone of cū are irregular.
\end{cousinweb}

\begin{grammarlens}[title={What you've built}]
One more word for describing things.
\end{grammarlens}

\subsection*{Your turn}
\begin{itemize}
\raggedright
  \item \emph{Say it:} \emph{cūxīn}
  \item \emph{Say it:} \emph{cūxīn}, once more
  \item \emph{Say it:} \emph{nǐ tài cūxīn le} — You're so careless.
  \item \emph{Recall:} say the Mandarin for hospitable, then the Mandarin for putting your heart into it, then say \emph{cūxīn} again
\end{itemize}

\begin{tcolorbox}[breakable,colback=teal!4,colframe=teal!35!black,title={Before you move on}]
What does cūxīn mean? (\textquotedblleft{}Careless\textquotedblright{}.) Say it once more.
\end{tcolorbox}
\section[xìxīn]{xìxīn — careful; attentive to detail}

\label{lesson:ZH-C273-xixin}



\begin{quote}
Before the new one: say the Mandarin for putting your heart into it, then the Mandarin for careless.
\end{quote}

\subsection*{You'll want to know: xìxīn}
\emph{xìxīn} — \textquotedblleft{}careful; attentive to detail\textquotedblright{}.

\begin{quote}
\emph{tā gōngzuò hěn xìxīn} — She's very careful in her work.
\end{quote}

Its characters come later; for now it is a word for the ear.

\begin{cousinweb}
\emph{xìxīn} is built from \emph{xì} 'fine' and \zh{心} 'the heart', both inherited from Middle Chinese, \emph{sejH} and \emph{sim} (Guangyun, in Baxter's transcription), which glosses the first 'small'.
\end{cousinweb}

\begin{grammarlens}[title={What you've built}]
One more word for describing things.
\end{grammarlens}

\subsection*{Your turn}
\begin{itemize}
\raggedright
  \item \emph{Say it:} \emph{xìxīn}
  \item \emph{Say it:} \emph{xìxīn}, once more
  \item \emph{Say it:} \emph{tā gōngzuò hěn xìxīn} — She's very careful in her work.
  \item \emph{Recall:} say the Mandarin for putting your heart into it, then the Mandarin for careless, then say \emph{xìxīn} again
\end{itemize}

\begin{tcolorbox}[breakable,colback=teal!4,colframe=teal!35!black,title={Before you move on}]
What does xìxīn mean? (\textquotedblleft{}Careful; attentive to detail\textquotedblright{}.) Say it once more.
\end{tcolorbox}
